\documentclass[10pt,letterpaper]{article}
\usepackage[T1]{fontenc}
\usepackage[utf8]{inputenc}
\usepackage{lmodern}
\usepackage[margin=0.76in]{geometry}
\usepackage[table]{xcolor}
\usepackage{array,tabularx,longtable,booktabs}
\usepackage{enumitem}
\usepackage{amssymb,amsmath}
\usepackage{hyperref}
\usepackage{fancyhdr,lastpage}
\usepackage{microtype}
\usepackage{parskip}
\definecolor{cdblue}{HTML}{203957}
\definecolor{cdgray}{HTML}{586677}
\hypersetup{colorlinks=true,linkcolor=cdblue,urlcolor=cdblue,pdftitle={Chapter 13: Managing Components and Dependencies}}
\pagestyle{fancy}
\fancyhf{}
\lhead{\small\color{cdgray}Continuous Delivery Engineering}
\rhead{\small\color{cdgray}Study guide}
\cfoot{\thepage\ / \pageref{LastPage}}
\setlength{\headheight}{14pt}
\setlength{\emergencystretch}{2em}
\setlist[itemize]{leftmargin=*,itemsep=3pt,topsep=4pt}
\setlist[enumerate]{leftmargin=*,itemsep=3pt,topsep=4pt}
\begin{document}
\begin{center}
{\LARGE\bfseries\color{cdblue}Chapter 13: Managing Components and Dependencies}\\[7pt]
{\large III: The Delivery Ecosystem}\\[6pt]
{\small Reference edition: 2026-10-30 $\bullet$ Original educational and operational material}
\end{center}
\noindent\color{cdblue}\rule{\linewidth}{1.2pt}\color{black}
\vspace{0.4em}
\section{Chapter scope and provenance}
Primary reading: Humble and Farley, Continuous Delivery (2010), Chapter 13, "Managing Components and Dependencies"; Part III: The Delivery Ecosystem. This guide is an original synthesis and not a reproduction of the chapter. The book provides the foundational delivery principles; current AWS, Terraform and GitHub Actions references in the applied sections are supplementary.


\section{Learning outcomes}
\begin{enumerate}
\item Identify dependency cycles and contract compatibility risks.
\item Use feature toggles and branch by abstraction.
\item Build component promotion rules.
\end{enumerate}

\section{Core conceptual model}
Focus areas: releasable changes; dependency graphs; binary repositories; abstraction.

\begin{itemize}
\item Input: a version-controlled change, explicit requirements, and a reproducible execution context.
\item Mechanism: controlled feedback, automation, evidence capture, and a stop condition where risk exceeds tolerance.
\item Output: a verifiable release candidate or a diagnosed failure with an accountable owner.
\end{itemize}

\section{Reading and analysis protocol}
\begin{enumerate}
\item Read the chapter and record five propositions, distinguishing an author claim from an example or assumption.
\item Identify two failure modes, one recovery technique, and one measurable acceptance condition.
\item Trace each technique to the stages and controls in the deployment-pipeline-architecture reference.
\item State how the proposed practice changes for a regulated system with segregated duties.
\end{enumerate}

\section{Detailed technical study notes}
\subsection{Small releasable increments}
New behavior can be merged behind toggles or abstractions, reducing long-lived branches and concurrent integration debt.

\subsection{Dependency contracts}
Producer and consumer version constraints must be explicit and tested; semantic version labels alone do not guarantee compatibility.

\subsection{Component packaging}
Split code into independently meaningful components only when their boundaries justify extra pipeline and provenance costs.

\subsection{Dependency graphs}
Detect transitive edges and cycles. Changes to shared interfaces trigger targeted downstream integration checks.

\subsection{Binary repositories}
Record artifact digest, publication permissions, retention and revocation; preserve previous known-good versions for recovery.

\subsection{Branch by abstraction}
Introduce a stable interface, migrate users incrementally and remove the old branch of behavior only after proving parity.

\section{Domain-specific challenge questions}
\begin{enumerate}
\item For small releasable increments, identify the precondition, expected evidence, adversarial failure case and measurable corrective action.
\item For dependency contracts, identify the precondition, expected evidence, adversarial failure case and measurable corrective action.
\item For component packaging, identify the precondition, expected evidence, adversarial failure case and measurable corrective action.
\item For dependency graphs, identify the precondition, expected evidence, adversarial failure case and measurable corrective action.
\end{enumerate}
\section{Guided exercise}
Create a dependency graph and classify which changes trigger downstream tests.

\subsection{Required evidence}
\begin{itemize}
\item A version-controlled design or executable artifact.
\item A failure demonstration, negative test or counterexample.
\item An improvement statement with baseline, target, owner and verification method.
\end{itemize}


\section{Engineering scenario}
A patched shared library requires ten teams to rebuild and redeploy simultaneously.

\begin{itemize}
\item Describe the failure propagation path and the earliest feasible detection point.
\item Recommend a minimal remediation and its rollback or containment plan.
\item Define a measure that would demonstrate risk reduction.
\end{itemize}

\section{Knowledge checks}
\begin{enumerate}
\item Which condition is the gate, and which action is the response when it fails?
\item What evidence distinguishes successful automation from an assumed successful release?
\item Where might environments, data, dependencies or people invalidate the technique?
\item What feedback could be moved earlier without losing necessary confidence?
\item Which assumptions are book principles, and which are modern supplementary implementation choices?
\end{enumerate}

\section{Completion rubric}
Score each dimension from 0 (missing) to 4 (demonstrated and independently verifiable): conceptual accuracy; reproducibility; negative testing; security and safety; documentation and traceability. Pass threshold: 16/20 with no zero on reproducibility or safety.


\section{Prerequisites and onward route}
Chapters 1--4 and prior dependency stages. Next: Chapters 14 onward.



\end{document}
